\section{Comparison with random transpositions}\label{cmp:section}
This section proves $H\succeq\tfrac12L_T$ and the resulting bounds on the
isotypic components. It uses the balance
\eqref{eq:balance} and the scalar \eqref{eq:content}; no estimate on the
Johnson scheme is needed. Section~\ref{root:section} consumes
Proposition~\ref{cmp:sectors}.

\begin{lemma}[The exchange is a trade]\label{cmp:swap}
Let $a=\{x,w\}$ and let $\tau$ be the transposition of $x$ and $w$. For every
$G\in\Omega$, the graph $\tau G$ lies in the $a$-fiber of $G$.
\end{lemma}

\begin{proof}
Relabeling by $\tau$ fixes every pair disjoint from $a$ and the pair $a$
itself. For $y\notin a$ it exchanges the pairs $xy$ and $wy$. So $\tau G$
agrees with $G$ at every common vertex and at every vertex adjacent to
neither endpoint, and the singletons adjacent to $x$ in $\tau G$ are those
adjacent to $w$ in $G$. Thus $\tau G$ agrees with $G$ off the pairs joining
$a$ to $S_a(G)$, and $\tau G\in\Omega$. This is the definition of the
$a$-fiber. In the description of Section~\ref{sec:setup}, $\tau$ replaces the
set $X$ by its complement in $S_a(G)$, which has $q$ elements by
\eqref{eq:balance}.
\end{proof}

\begin{proposition}[Transposition comparison]\label{cmp:main}
For every pair $a$, put $Q_a=\tfrac12(I+T_a)$. Then $Q_a-E_a$ is an orthogonal
projection,
\begin{equation}\label{cmp:local}
 h_a=\tfrac12(I-T_a)+(Q_a-E_a)\succeq\tfrac12(I-T_a),
\end{equation}
and $H=\tfrac12L_T+R$, where $R=\sum_a(Q_a-E_a)\succeq0$ commutes with every
$T_\sigma$.
\end{proposition}

\begin{proof}
The operator $Q_a$ is the orthogonal projection onto the functions fixed by
$T_a$, since $T_a$ is an orthogonal involution. A function in the range of
$E_a$ is constant on every $a$-fiber. By Lemma~\ref{cmp:swap} it takes the
same value at $G$ and $\tau G$, so it is fixed by $T_a$. Hence the range of
$E_a$ is contained in the range of $Q_a$. For orthogonal projections this
gives $Q_aE_a=E_aQ_a=E_a$, so $Q_a-E_a$ is the orthogonal projection onto the
orthogonal complement of the range of $E_a$ inside the range of $Q_a$.
Then $h_a=(I-Q_a)+(Q_a-E_a)$ and $I-Q_a=\tfrac12(I-T_a)$, which is
\eqref{cmp:local}. Summing over $a$ gives $H=\tfrac12L_T+R$. Both $H$ and
$L_T$ commute with every $T_\sigma$, by \eqref{eq:equivariance} and because
conjugation permutes the transpositions; so does $R$.
\end{proof}

\begin{lemma}[Content bound]\label{cmp:content}
Let $\lambda$ be a partition of $n\ge4$ with $\lambda_1\le n-2$. Then
\[
 \operatorname{cont}(\lambda)\le\binom{n-2}2-1,\qquad
 \binom n2-\operatorname{cont}(\lambda)\ge2n-2,
\]
with equality in both for $\lambda=(n-2,2)$.
\end{lemma}

\begin{proof}
Put $\lambda_1=n-k$, so that $2\le k\le n-1$ and the rows $i\ge2$ contain $k$
cells in total. Since $\binom u2+\binom v2\le\binom{u+v}2$ for nonnegative
integers and $i-1\ge1$ in these rows,
\[
 \operatorname{cont}(\lambda)
 =\binom{n-k}2+\sum_{i\ge2}\Bigl[\binom{\lambda_i}2-(i-1)\lambda_i\Bigr]
 \le\binom{n-k}2+\binom k2-k=:\phi(k).
\]
The function $\phi$ is a convex quadratic polynomial in $k$, and
\[
 \phi(2)=\binom{n-2}2-1,\qquad
 \phi(n-1)=\binom{n-1}2-(n-1)=\frac{(n-1)(n-4)}2=\binom{n-2}2-1 .
\]
So $\phi(k)\le\binom{n-2}2-1$ for $2\le k\le n-1$. The second inequality
follows from $\binom n2-\binom{n-2}2=2n-3$. For $\lambda=(n-2,2)$ the contents
of the second row are $-1$ and $0$, so
$\operatorname{cont}(\lambda)=\binom{n-2}2-1$.
\end{proof}

\begin{proposition}[Isotypic components]\label{cmp:sectors}
For every partition $\lambda$ of $n$,
\[
 H\succeq\frac12\Bigl[\binom n2-\operatorname{cont}(\lambda)\Bigr]I
 \quad\text{on }U_\lambda .
\]
The right side is at least $n-1$ when $\lambda\notin\{(n),(n-1,1)\}$, equals
$n/2$ for $\lambda=(n-1,1)$, and equals $0$ for $\lambda=(n)$.
\end{proposition}

\begin{proof}
The first assertion follows from Proposition~\ref{cmp:main} and
\eqref{eq:content}, because $U_\lambda$ is invariant under $H$ and $L_T$. A
partition other than $(n)$ and $(n-1,1)$ has $\lambda_1\le n-2$, and
Lemma~\ref{cmp:content} applies. Finally
$\operatorname{cont}(n-1,1)=\binom{n-1}2-1$ and
$\operatorname{cont}(n)=\binom n2$.
\end{proof}

\begin{remark}[Unequal degrees]\label{cmp:general}
Lemma~\ref{cmp:swap} and \eqref{cmp:local} hold on $\Omega_d$ for an arbitrary
graphical degree sequence $d$ and every pair $a=\{x,w\}$ with $d_x=d_w$, with
the same proofs: equal degrees give \eqref{eq:balance} for that pair, and the
transposition of $x$ and $w$ preserves $\Omega_d$. Pairs of unequal degree
admit no such comparison, since the exchange does not preserve $\Omega_d$.
\end{remark}
